% BEGIN_REPORT::WRITE::WRT-0002::2026-10-07T16:47:38-07:00
\documentclass[11pt]{article}
\usepackage[margin=1in]{geometry}
\usepackage{amsmath,amssymb}
\usepackage{xurl,hyperref}
\title{WRT-0002: Zero-Child VALUE Members Are Exact Partial-Portfolio Plans}
\author{WRITE, ARC/WRAPPER research team}
\date{7 October 2026}
\begin{document}
\maketitle

\section{Result and strict status}
\textbf{ZERO-CHILD-PORTFOLIO1 is PROVED/EXACT.}  It is a specialization of
CERTIFIED-PARTIAL-PORTFOLIO1: a fixed portfolio member may return a
self-certified constant combiner with $q=0$ children.  This does not evade
the theorem's exact-success condition: the member must establish the answer
before it reports SUCCESS.  The positive-cofactor member recovered in
RED-0002 meets that condition on its declared domain.  It does not cover
arbitrary SAT instances, so the joint WRAPPER interface, the SAT-specific ARC
successor, and P versus NP remain open.

This run also closes a procedural stall: RED-0002's GitHub bodies had been
verified, while its local log and outbox manifest still said RUNNING.  Their
state is repaired to VERIFIED using the recorded commit and blob identifiers.

\section{Claim card}
\begin{description}
\item[] \textbf{Claim ID:} ZERO-CHILD-PORTFOLIO1.
\item[] \textbf{Exact quantifiers:} fix a finite family of uniform partial
procedures.  Every member halts with FAIL or an exact finite plan; all
attempts are charged; and every input is covered by some successful member.
The case $q_j=0$ is allowed, with a constant Boolean combiner.  Sequential
first-SUCCESS selection remains exact.
\item[] \textbf{Model and measure:} ordinary indexed CNF, forward-only plans,
charged generation and live state, and ordinary bit time/space.  For a
zero-child plan, child generation, access, and recursive acquisition are
vacuous, but the successful VALUE certificate and its input scan remain
charged.
\item[] \textbf{SAME\_REALIZATION object:} the member's SUCCESS record and
constant combiner produced by the same run that performs the certified scan.
\item[] \textbf{Weakest useful theorem:} make the $q=0$ specialization
explicit without claiming that any finite collection of partial members
covers general SAT.
\item[] \textbf{Falsifier:} an input on which a member reports a constant
SUCCESS value different from SAT, or a fixed portfolio claimed to be total
while every member returns FAIL.
\end{description}

\section{Corpus and source review}
The run read the proof-spines protocol, \texttt{AGENTS.md}, \texttt{CODEX.md},
the repository research skill, root \texttt{main.tex} and
\texttt{main2.tex}--\texttt{main5.tex}, the live peer reports, raw outputs,
and the implementation for RED-0002.  GitHub readback confirmed the report,
code, raw result, bibliography, and both spines.  Drive search then opened the
underlying historical backdoor-cost, $q$-Horn portfolio, recursive
normalization, streaming-composition, and orbit-batching documents rather
than relying on their labels.  They supply conditional interfaces or
restricted constructions; none proves the missing favorable general-SAT
coverage premise.

The original backdoor-DNF paper defines the residuals and establishes the
tautological cover in its Theorem~2; its Theorem~3 is a parameterized
detection result for the four stated base-class choices
\cite{maehlmann2021}.  Neither theorem supplies the desired joint WRAPPER
resource bound.  For a residual already known to be Horn, Dowling--Gallier
supplies linear-time recognition/solving in its ordinary explicit model
\cite{dowling1984horn}.  The transfer used here is narrower: the displayed
positive-cofactor condition can certify TRUE directly, without materializing
or solving any residual.  These are theorem transfers with their stated
hypotheses, not analogies to an arbitrary semantic compressor.

\section{Proof of the specialization}
In CERTIFIED-PARTIAL-PORTFOLIO1, a successful plan is
$(C_j;G_{j,1},\ldots,G_{j,q_j})$ and must satisfy
\[
 \operatorname{SAT}(F)=C_j(\operatorname{SAT}(G_{j,1}),\ldots,
 \operatorname{SAT}(G_{j,q_j})).
\]
Put $q_j=0$.  A Boolean function of zero arguments is a constant, so exactness
is precisely $\operatorname{SAT}(F)=C_j()$.  The fanout and child-size
conditions are vacuous and may be recorded with $a_j=b_j=0$.  The theorem's
success/failure time and space premise still charges the scan that justifies
the constant answer.  The existing proof then applies unchanged: coverage
selects some exact successful member, the finite sequence of failed attempts
has maximum rather than summed exponent, and failed state is discarded.

For RED-0002's positive-cofactor member, fix $p$.  It scans every clause.  If
a clause contains $\neg p$, then setting $p=0$ satisfies it.  Otherwise the
member requires the clause, after deleting literals on $p$, to be a nonempty
positive clause.  Setting $p=0$ and all remaining variables to TRUE therefore
satisfies every clause.  Only after verifying this syntactic condition does
the member return the zero-child constant TRUE plan.  Failure is harmless;
claiming TRUE without the scan is not.

\section{Reproducible bounded test}
The command
\begin{verbatim}
timeout 30s python3 experiments/proof_spines/
  zero_child_portfolio_test.py
\end{verbatim}
ran under an internal 25-second alarm.  It exhaustively checked all 12,616
normalized width-two formulas on three variables with at most five clauses.
The positive-cofactor member returned 638 exact zero-child TRUE plans; the
fixed positive-cofactor/Horn/2-CNF portfolio agreed with exhaustive truth
tables on all formulas, with first-success counts 638, 4,824, and 7,154.
There were zero violations.  The code SHA256 is
\path{905bfaa448e6fa741a78f8226d0967922ce89e86ef99831ecb331b9163696f93}.
Raw output is
\path{results/raw/proof_spines/2026-10-07_164738_zero_child_portfolio.json}.
The computation checks the implementation; the argument above supplies the
general theorem specialization.

\section{Hostile controls and limitations}
\begin{description}
\item[] \textbf{Unsafe first finisher:} constant TRUE fails on
$(p)\wedge(\neg p)$, while the certified member returns FAIL.
\item[] \textbf{Easy family:} source-paper instances through $s=10$ all
produce zero children.  This establishes only the declared success domain.
\item[] \textbf{Renaming and padding:} after swapping the distinguished
variable, the fixed-name member may safely fail.  Thirty-two tautological
padding clauses normalize away; padding yields no asymptotic promotion.
\item[] \textbf{Independent products:} four independently accepted
components combine by AND without child generation.  This is a semantic
closure check, not a worst-case acquisition theorem.
\item[] \textbf{Source carry and current-scale calls:} the member uses no SAT
call and carries no witness/source map.  It charges its scan but does not
solve instances outside its syntactic domain.
\item[] \textbf{Oracle and algebraic replacement:} the specialization
relativizes and supplies no lower-bound or algebrization hinge.
\item[] \textbf{Finite versus growing depth:} this is absorption into one
fixed finite portfolio.  It proves neither a growing-depth uniform algorithm
nor a cofinal ARC return map.
\end{description}

\section{Manuscript reconciliation and handoff}
Root \texttt{main5.tex} now states the $q=0$ specialization, removes stale
``conditional pending review'' language, and records the independent test.
Root \texttt{main4.tex} records that the result has no ARC mathematical
status effect.  No citation change is needed because the verified primary
sources are already in \texttt{spines.bib}.

\begin{description}
\item[] \textbf{Handoff to LEAD/PROOF:} retain zero-child VALUE members in the
fixed partial portfolio, but target a coverage theorem on arbitrary original-
scale SAT instances.  The next candidate must jointly charge the classifier
or acquisition routine and show that some favorable member succeeds on every
input; merely adding more easy-family VALUE members cannot hit a cofinal ARC
seed.
\end{description}

\bibliographystyle{plain}
\bibliography{spines}
\end{document}
% END_OF_VERIFIED_REPORT::WRITE::WRT-0002
